% GENERATED FILE. Edit canonical lessons, then run npm run generate:books.
% canonical-source-hash: fnv1a64:f6c8fe3e105b24e7
% canonical-lessons: RU-C454-rasprodazha, RU-C454-utsenka, RU-C454-aktsiya, RU-C454-torgovatsya, RU-C454-pereplachivat

\chapter{Sales and bargains}
\label{ch:ru-sales-and-bargains}
\hlchaptermodality{454}

\begin{chapteropening}
\textbf{By the end of this chapter:} \emph{I can find a bargain: the sale, a markdown, a special offer, haggling and paying too much.}

Complete the last lesson of chapter 454: I can find a bargain: the sale, a markdown, a special offer, haggling and paying too much.
\end{chapteropening}

\section[\texorpdfstring{\ru{распродажа} (rasprodázha)}{распродажа (rasprodázha)}]{\ru{распродажа} (rasprodázha) — a sale (clearance)}

\label{lesson:RU-C454-rasprodazha}



\begin{quote}
Before the new one: say the Russian for to pinch, then the Russian for to exchange.
\end{quote}

\subsection*{You'll want to know: \ru{распродажа}}
\textbf{\ru{распродажа}} (\emph{rasprodázha}) — \textquotedblleft{}a sale (clearance)\textquotedblright{}.

\begin{quote}
\ru{В} \ru{январе} \ru{во} \ru{всех} \ru{магазинах} \ru{распродажа}.
\end{quote}

\begin{quote}
\emph{V yanvaré vo vsekh magazínakh rasprodázha.}
\end{quote}

\begin{quote}
In January there is a sale in all the shops.
\end{quote}

Feminine. \emph{\ru{На} \ru{распродаже}} — in the sale.

\begin{cousinweb}
\emph{\ru{распродажа}} is built from \emph{\ru{рас}-} 'off, all over' and \emph{\ru{продажа}} 'a sale', from \emph{\ru{продать}} 'to sell', on \emph{\ru{дать}} 'to give'. EtymDB lists \emph{\ru{продажа}} without a deeper source.
\end{cousinweb}

\begin{grammarlens}[title={What you've built}]
One more word for telling, explaining and describing at B1.
\end{grammarlens}

\subsection*{Your turn}
\begin{itemize}
\raggedright
  \item \emph{Say it:} \emph{rasprodázha}
  \item \emph{Say it:} \emph{rasprodázha}, once more
  \item \emph{Say it:} a sentence of your own with \emph{rasprodázha}
  \item \emph{Recall:} say the Russian for to pinch, then the Russian for to exchange, then say \emph{rasprodázha} again
\end{itemize}

\begin{tcolorbox}[breakable,colback=teal!4,colframe=teal!35!black,title={Before you move on}]
What does \ru{распродажа} mean? (\textquotedblleft{}A sale (clearance)\textquotedblright{}.) Say it once more.
\end{tcolorbox}
\section[\texorpdfstring{\ru{уценка} (utsénka)}{уценка (utsénka)}]{\ru{уценка} (utsénka) — a markdown, a price reduction}

\label{lesson:RU-C454-utsenka}



\begin{quote}
Before the new one: say the Russian for to exchange, then the Russian for a sale.
\end{quote}

\subsection*{You'll want to know: \ru{уценка}}
\textbf{\ru{уценка}} (\emph{utsénka}) — \textquotedblleft{}a markdown, a price reduction\textquotedblright{}.

\begin{quote}
\ru{На} \ru{эту} \ru{куртку} \ru{большая} \ru{уценка}.
\end{quote}

\begin{quote}
\emph{Na étu kúrtku bolsháya utsénka.}
\end{quote}

\begin{quote}
This jacket has a big markdown.
\end{quote}

Feminine. \emph{\ru{Товар} \ru{по} \ru{уценке}} — reduced goods.

\begin{cousinweb}
\emph{\ru{уценка}} is built from \emph{\ru{у}-} 'off' and \emph{\ru{цена}} 'a price', which is inherited from Proto-Slavic \emph{cěna} (EtymDB).
\end{cousinweb}

\begin{grammarlens}[title={What you've built}]
One more word for telling, explaining and describing at B1.
\end{grammarlens}

\subsection*{Your turn}
\begin{itemize}
\raggedright
  \item \emph{Say it:} \emph{utsénka}
  \item \emph{Say it:} \emph{utsénka}, once more
  \item \emph{Say it:} a sentence of your own with \emph{utsénka}
  \item \emph{Recall:} say the Russian for to exchange, then the Russian for a sale, then say \emph{utsénka} again
\end{itemize}

\begin{tcolorbox}[breakable,colback=teal!4,colframe=teal!35!black,title={Before you move on}]
What does \ru{уценка} mean? (\textquotedblleft{}A markdown, a price reduction\textquotedblright{}.) Say it once more.
\end{tcolorbox}
\section[\texorpdfstring{\ru{акция} (áktsiya)}{акция (áktsiya)}]{\ru{акция} (áktsiya) — a special offer, a promotion; a share (in a company)}

\label{lesson:RU-C454-aktsiya}



\begin{quote}
Before the new one: say the Russian for a sale, then the Russian for a markdown.
\end{quote}

\subsection*{You'll want to know: \ru{акция}}
\textbf{\ru{акция}} (\emph{áktsiya}) — \textquotedblleft{}a special offer, a promotion; a share (in a company)\textquotedblright{}.

\begin{quote}
\ru{По} \ru{акции} \ru{второй} \ru{кофе} — \ru{бесплатно}.
\end{quote}

\begin{quote}
\emph{Po áktsii vtoróy kófe — besplátno.}
\end{quote}

\begin{quote}
On special offer, the second coffee is free.
\end{quote}

Feminine. \emph{\ru{По} \ru{акции}} — on special offer. Also a share on the stock market, or a protest.

\begin{cousinweb}
\emph{\ru{акция}} is usually traced, probably through German or French, to Latin \emph{āctiō} 'a doing, an action' (EtymDB lists the Latin word), from \emph{agere} 'to drive, to do'.
\end{cousinweb}

\begin{grammarlens}[title={What you've built}]
One more word for telling, explaining and describing at B1.
\end{grammarlens}

\subsection*{Your turn}
\begin{itemize}
\raggedright
  \item \emph{Say it:} \emph{áktsiya}
  \item \emph{Say it:} \emph{áktsiya}, once more
  \item \emph{Say it:} a sentence of your own with \emph{áktsiya}
  \item \emph{Recall:} say the Russian for a sale, then the Russian for a markdown, then say \emph{áktsiya} again
\end{itemize}

\begin{tcolorbox}[breakable,colback=teal!4,colframe=teal!35!black,title={Before you move on}]
What does \ru{акция} mean? (\textquotedblleft{}A special offer, a promotion; a share (in a company)\textquotedblright{}.) Say it once more.
\end{tcolorbox}
\section[\texorpdfstring{\ru{торговаться} (torgovátsya)}{торговаться (torgovátsya)}]{\ru{торговаться} (torgovátsya) — to haggle, to bargain}

\label{lesson:RU-C454-torgovatsya}



\begin{quote}
Before the new one: say the Russian for a markdown, then the Russian for a special offer.
\end{quote}

\subsection*{You'll want to know: \ru{торговаться}}
\textbf{\ru{торговаться}} (\emph{torgovátsya}) — \textquotedblleft{}to haggle, to bargain\textquotedblright{}.

\begin{quote}
\ru{На} \ru{рынке} \ru{принято} \ru{торговаться}.
\end{quote}

\begin{quote}
\emph{Na rýnke prínyato torgovátsya.}
\end{quote}

\begin{quote}
At the market it is usual to haggle.
\end{quote}

Imperfective; the perfective is \emph{\ru{сторговаться}}. Present \emph{\ru{торгуюсь}, \ru{торгуешься}}.

\begin{cousinweb}
\emph{\ru{торговаться}} is built on \emph{\ru{торг}} 'trade, a market', which is inherited through Old East Slavic \emph{tŭrgŭ} from Proto-Slavic \emph{tŭrgŭ} (EtymDB).
\end{cousinweb}

\begin{grammarlens}[title={What you've built}]
One more word for telling, explaining and describing at B1.
\end{grammarlens}

\subsection*{Your turn}
\begin{itemize}
\raggedright
  \item \emph{Say it:} \emph{torgovátsya}
  \item \emph{Say it:} \emph{torgovátsya}, once more
  \item \emph{Say it:} a sentence of your own with \emph{torgovátsya}
  \item \emph{Recall:} say the Russian for a markdown, then the Russian for a special offer, then say \emph{torgovátsya} again
\end{itemize}

\begin{tcolorbox}[breakable,colback=teal!4,colframe=teal!35!black,title={Before you move on}]
What does \ru{торговаться} mean? (\textquotedblleft{}To haggle, to bargain\textquotedblright{}.) Say it once more.
\end{tcolorbox}
\section[\texorpdfstring{\ru{переплачивать} (perepláchivat)}{переплачивать (perepláchivat)}]{\ru{переплачивать} (perepláchivat) — to overpay, to pay too much}

\label{lesson:RU-C454-pereplachivat}



\begin{quote}
Before the new one: say the Russian for a special offer, then the Russian for to haggle.
\end{quote}

\subsection*{You'll want to know: \ru{переплачивать}}
\textbf{\ru{переплачивать}} (\emph{perepláchivat}) — \textquotedblleft{}to overpay, to pay too much\textquotedblright{}.

\begin{quote}
\ru{В} \ru{туристских} \ru{местах} \ru{мы} \ru{всегда} \ru{переплачиваем}.
\end{quote}

\begin{quote}
\emph{V turístskikh mestákh my vsegdá perepláchivayem.}
\end{quote}

\begin{quote}
In touristy places we always pay too much.
\end{quote}

Imperfective; perfective \emph{\ru{переплатить}} (\emph{\ru{переплачу}, \ru{переплатишь}}). Takes \emph{\ru{за}} + accusative.

\begin{cousinweb}
\emph{\ru{переплачивать}} is built from \emph{\ru{пере}-} 'over, too much' and \emph{-\ru{плачивать}}, from \emph{\ru{платить}} 'to pay', which EtymDB takes from Proto-Slavic \emph{platiti}.
\end{cousinweb}

\begin{grammarlens}[title={What you've built}]
One more word for telling, explaining and describing at B1.
\end{grammarlens}

\subsection*{Your turn}
\begin{itemize}
\raggedright
  \item \emph{Say it:} \emph{perepláchivat}
  \item \emph{Say it:} \emph{perepláchivat}, once more
  \item \emph{Say it:} a sentence of your own with \emph{perepláchivat}
  \item \emph{Recall:} say the Russian for a special offer, then the Russian for to haggle, then say \emph{perepláchivat} again
\end{itemize}

\begin{tcolorbox}[breakable,colback=teal!4,colframe=teal!35!black,title={Before you move on}]
What does \ru{переплачивать} mean? (\textquotedblleft{}To overpay, to pay too much\textquotedblright{}.) Say it once more.
\end{tcolorbox}
